\lecture{1}{Thu 25 May 2023 02:21}{Summer Research}

\section{Proof of convergence of AF-SIRW to BMPE when $p < \frac{1}{2}$}

The \textit{encountered total drift} at a point $x$ is defined as
\[
\delta_{x,m} = \sum_{i=1}^{\infty} \mathbb{E}\left[ X_{i+1} - X_i | \mathcal{F}_i^X \right] 
\mathbf{1}\left( X_i = x \right) 
\mathbf{1}\left(\mathcal{L}_{x,i} \le m\right)
.\] 

Consider $y > x$.

We define 
 \begin{align*}
	 \Delta_y^{(x,m)} &= \sum_{i=1}^{\mathcal{L}_y^{(x,m)}} \mathbb{E}\left[ X_{i+1} - X_i | \mathcal{F}_i^X \right] \mathbf{1}(X_i = y)\\
	 \rho_y^{(x,m)} &= \mathbb{E}\left[ \Delta_y^{(x,m)} | \mathcal{G}_{y-1} ^{(x,m)} \right] 
.\end{align*} 

where $\mathcal{G}_{z}^{(x,m)} := \sigma\left( \mathcal{E}_y^{(x,m)}: x \le y \le z \right) $

We will decompose the sum $\delta_{x,m}$ into a martingale difference and a compensation term. First we would like to establish

\[
	\sum_{y = x+1}^{\mathcal{M}_{x,m}} \Delta_y^{(x,m)} - \rho_y^{(x,m)} = o(\sqrt{n}) 
.\] 

This is established by a variance bound
\[
	P\left( \left| \sum_{y} \Delta_y^{(x,m)} - \rho_y^{(x,m)} \right| > \varepsilon\sqrt{n}  \right)
	\le  \frac{\operatorname{Var} \left( \sum_{y} \Delta_y^{(x,m)} - \rho_y^{(x,m)}\right) }{\varepsilon^2 n}
	\to 0
.\] 

Using tightness lemma, there are $O(\sqrt{n})$ terms in the sum.
To achieve desired bound, we would expect for any $y$, $\operatorname{Var}\left[ \Delta_y^{(x,m)} - \rho_y^{(x,m)}\right] = o(\sqrt{n})$. (Suppose, that covariance is negligible here).

The second part is to establish
\[
	\sum_{y = x+1}^{\mathcal{M}_{x,m}} \rho_y^{(x,m)} - \rho = o(\sqrt{n}) 
.\] 

for some fixed $\rho$. We expect that $\rho = \gamma$ as in the paper, but the convergence $\rho_y^{(x,m)} \to \rho$ is less straightforward to establish.

\subsection{Heuristic $O(1)$ bound for local variance}

\textbf{Notations.} For simplicity of notation, define $\alpha(i,j) := \frac{- l(i) + r(j)}{l(i) + r(j)} = \frac{-w(2 i - 1) + w(2j)}{w(2i-1)+w(2j)}$. This is the encountered drift at some node when the Polya Urn state is $(i,j)$.

Write $\xi_i$ for the number of descendents of the $i$-th ancestor in the BLP. Then we have the branching process representation of $\rho_y^{(x,m)}$ :
\[
	\rho_y^{(x,m)} = \mathbb{E}\left[ \sum_{i = 1}^{\mathcal{E}_{y-1}^{(x,m)}} \sum_{j = 0}^{\xi_i} \alpha(i, j + \sum_{k < i} \xi_k) \right] 
.\] 

The heuristic computation is as follows.
 
First compute variance by conditioning.
 \begin{align*}
	  &\operatorname{Var} \left[ \Delta_y^{(x,m)} - \rho_y^{(x,m)}  \right]\\
	 &=  \mathbb{E} \left[\operatorname{Var} \left[ \Delta_y^{(x,m)} - \rho_y^{(x,m)} | \mathcal{E}_{y-1}^{(x,m) } = n \right]\right]\\
	 &= \mathbb{E}\left[ \operatorname{Var} \left[ \sum_{i = 1}^n \sum_{j = 0}^{\xi_i} \alpha(i, j + \sum_{< i} \xi_.) \right]  \right] 
.\end{align*}

Now begin the heuristic calculation. Consider very large $n$ and sum from $\frac{n}{2}$ to $n$. In this way we may assume that all $\xi_i$ are arbitrarily close to a $\text{Geom}(\frac{1}{2})$ random variable.

\begin{align*}
	&\operatorname{Var} \left[ \sum_{i = \frac{n}{2}}^n \sum_{j = 0}^{\xi_i} \alpha(i, j + \sum_{< i} \xi_.) \right]\\
	&\approx \sum_{i = \frac{n}{2}}^n \operatorname{Var} \left[  \sum_{j = 0}^{\xi_i} \alpha(i, j + \sum_{< i} \xi_.) \right] \intertext{where we assume the covariances are small.} \\
	&\approx \sum_{i = \frac{n}{2}}^n \operatorname{Var} \left[  \xi_i \alpha(i, \sum_{\le  i} \xi_.) \right] \intertext{since $\alpha$ is not sensitive to small perturbation to the second argument.} \\
	&\approx \sum_{i = \frac{n}{2}}^n \operatorname{Var} \left[  \alpha(i, \sum_{\le  i} \xi_.) \right] + 2 \mathbb{E}\left[ \left( \alpha(i, \sum_{\le i} \xi_.) \right) ^2 \right] \intertext{by assuming $\xi_i \sim \text{Geom}(\frac{1}{2})$.} 
\end{align*}

Approximate $\sum_{\le i} \xi_.$ by $N(i, 2i)$, we may write out an integral for $\operatorname{Var} \left[\alpha(i, N(i, 2i))\right]$ and obtain the order of magnitude estimate
\[
	\operatorname{Var} \left[\alpha(i, \sum_{\le i}\xi_.)\right] = O(i^{-2p-1})
.\] 

Summing over $i$, we obtain
\[
	\sum_{i = \frac{n}{2}}^n \operatorname{Var} \left[  \alpha(i, \sum_{\le  i} \xi_.) \right] = O(n^{-2p})
.\] 

By repeatedly apply the bound, halfing $n$ each time until some fixed constant $N$, we obtain
\[
	\sum_{i = N}^n \operatorname{Var} \left[  \alpha(i, \sum_{\le  i} \xi_.) \right]  = O(N^{-2p}) + O((2N)^{-2p)} + \ldots + O\left( (\log_2 \frac{n}{N} N)^{-2p} \right) = O(1) 
.\] 

\begin{note}{Problems about this estimation}
	Confusing notation was used in this part; $n$ should refer to the total number of steps taken by the walk, but it was also used to refer to the value of $\mathcal{E}_y^{(x, m)}$. These two things have different scale, in fact, we have approximately $\mathcal{E}_y^{(x,m)} = O(\sqrt{n} )$. From this we see that the ``cheapest'' estimate of covariance using variances, actually works.

	The non-diagonal covariance terms turns out to be non-negligible, in fact of the same scale as the diagonal terms, and not dependent on the distance from diagonal.
\end{note}

\subsection{Justifications of the variance bound}

Our definition is related to the paper by
\[
	\sum_{\le i}\xi_. \equiv \tau_i^{\mathcal{B}} - i
.\] 

Then according to Lemma 4.1 and Remark 4.2,

\begin{lemma}
	\[
		P\left( \left|   \sum_{\le i}\xi_. - i\right| \ge m \right)  \le C e ^{\frac{-c m^2}{ m \vee n}}
	.\] 
\end{lemma}

\begin{lemma}
	\[
		\mathbb{E}\left[ \alpha\left(i, \tau_i^{\mathcal{B}} - i\right) \right] \le O(i^{-p-\frac{1}{2}})
	.\] 
\end{lemma}
\begin{proof}
	Calculate by splitting into good and bad events. Define good event $G = \{|\tau_i^{\mathcal{B}} - 2i| < m \sqrt{ i} \} $. 

\begin{align*}
	\mathbb{E}\left[ \alpha\left(i, \tau_i^{\mathcal{B}} - i\right) \right]
	&= \mathbb{E}\left[ \alpha\left(i, \tau_i^{\mathcal{B}} - i\right) \mathbf{1}(G)\right] + \mathbb{E}\left[ \alpha\left(i, \tau_i^{\mathcal{B}} - i\right) \mathbf{1}(G^c)\right] \\
	&\le \mathbb{E}\left[ \alpha\left(i, \tau_i^{\mathcal{B}} - i\right)  | G\right] + P(G^c) \\
	&= \mathbb{E}\left[ \alpha\left(i, \tau_i^{\mathcal{B}} - i\right)  | G\right] + O(e^{-c m^2}) 
.\end{align*}

After bounding possible $j$'s, follow the argument in proof of Lemma 2.3.
\begin{align*}
	 \mathbb{E}\left[ \alpha\left(i, \tau_i^{\mathcal{B}} - i\right)  | G\right]
	 &\le \sup _{|i - j| < m \sqrt{i} } \alpha(i,j) \\
	 &\le C \sup \left| 2^p B \left( \frac{1}{(2i - 1)^p}  - \frac{1}{(2j)^p} + O\left(\frac{1}{i^{\varkappa
	 + 1}}\right) \right) \right|  \\
	 &\le O\left(\frac{1}{i^{\varkappa + 1}}\right) 
	 + C \sup \left| \frac{1}{i^p} - \frac{1}{j^p}\right| + C \sup \left| \frac{1}{(2i-1)^p} - \frac{1}{(2i)^p} \right| \\
	 &\le O\left(\frac{1}{i^{\varkappa + 1}}\right) 
	 + C m \sqrt{ i} \left( i - m \sqrt{ i}  \right) ^{- p - 1} + C 2^p (2 i - 1)^{ -p - 1}\\
	 &= O\left(i^{- p - \frac{1}{2}}\right) 
.\end{align*}

\end{proof}

\begin{lemma}
	\[
		\operatorname{Var} \left[ \alpha(i, \tau_i^{\mathcal{B}} - i) \right] \le    O(i^{- 2 p - 1})
	.\] 
	\label{lemma:variance-diagonal}
\end{lemma}
\begin{proof}
	Since we have controlled expectation, it remains to control $\mathbb{E} [\alpha^2]$. Follow the same procedure, we obtain
\begin{align*}
	 \mathbb{E}\left[ \alpha\left(i, \tau_i^{\mathcal{B}} - i\right)^2  | G\right]
	 &\le \sup _{|i - j| < m \sqrt{i} } \alpha(i,j)^2 \\
	 &\le C \sup \left| 2^p B \left( \frac{1}{(2^i - 1)^p}  - \frac{1}{(2j)^p} + O(\frac{1}{i^{\varkappa
	 + 1}}) \right) \right|^2  \\
	 &= O(i^{-2p-1}) 
.\end{align*}
\end{proof}

\begin{lemma}
	There exists $\delta>0$, such that for $\frac{n}{2} \le i < I \le n$, whenever $i_2 - i_1 > n^\delta$, we have
	\[
		\operatorname{Cov} \left(
			\sum_{j = 0}^{\xi_{i}}  \alpha(i,\sum_{\le i}\xi_.)\,\,, \quad
			\sum_{j = 0}^{\xi_{I}}  \alpha(I,\sum_{\le I}\xi_.)
		\right) = o(n^{-2p-2})
	.\] 
\end{lemma}
note: This lemma is experimentally shown to be \textbf{false}.

\begin{proposition}
	If $X_1 \in \mathcal{F}_1$ and $X_2 \in \mathcal{F}_1 \cup \mathcal{F}_2$, where $\mathcal{F}_1, \mathcal{F}_2$ are independent, then
	\[
		\operatorname{Cov} (X_1, X_2) = \operatorname{Cov} \left(X_1, \mathbb{E}\left[ X_2 | \mathcal{F}_1 \right] \right)
	.\] 
\end{proposition}


\begin{proof}
	\begin{align*}
		\operatorname{Cov} (X_1, X_2)
		&= \mathbb{E}\left[ \mathbb{E}\left[ (X_1 - \mathbb{E} X_1) (X_2 - \mathbb{E} X_2) | \mathcal{F}_1 \right]  \right]  \\
		&= \mathbb{E}\left[(X_1 - \mathbb{E} X_1)  \mathbb{E}\left[ (X_2 - \mathbb{E} X_2) | \mathcal{F}_1 \right]  \right]  \\
		&= \mathbb{E}\left[(X_1 - \mathbb{E} X_1)  (\mathbb{E}\left[ X_2  | \mathcal{F}_1 \right]- \mathbb{E} X_2)  \right]  \\
		&= \mathbb{E}\left[(X_1 - \mathbb{E} X_1)  (\mathbb{E}\left[ X_2  | \mathcal{F}_1 \right]- \mathbb{E} \left[ \mathbb{E}\left[ X_2  | \mathcal{F}_1 \right] \right] )  \right]  \\
		&= \operatorname{Cov} (X_1, \mathbb{E}\left[ X_2 | \mathcal{F}_1 \right] )
	.\end{align*}
\end{proof}


